\subsubsection{Cross-Venue Routing and Fragmentation}

\begin{principle} \hlt{Routing After Fragmentation}\\
In a fragmented market, liquidity is no longer represented by one order book. The execution problem becomes a joint choice of price, size, time, venue, fees, and information exposure.
\end{principle}

\begin{definition} \hlt{Smart Order Routing}\\
A smart order router (SOR) maps a trading instruction to venues according to displayed quotes, depth, historical fill rates, queue position, fees, latency, and order-type constraints. The execution algorithm usually chooses the parent schedule; the router chooses the spatial allocation of each child order.
\end{definition}

\begin{principle} \hlt{Temporal Then Spatial Split}\\
Optimal execution first chooses how fast to trade, balancing market impact against timing risk. In a fragmented market, each child order then needs a spatial split across venues, where the SOR tries to obtain at least the limit price and finish the child order quickly.
\end{principle}

\begin{process} \hlt{SOR Feedback Loop}
\begin{enumerate}[label=\roman*.]
\setlength{\itemsep}{0pt}
\item observe venue quotes, depths, fees, and recent fills;
\item decide the venue split under the execution policy;
\item send orders through venue gateways;
\item receive execution reports and market-data updates;
\item update the next routing decision.
\end{enumerate}
The decision loop is bounded by CPU time, network latency, matching-engine latency, and feed latency.
\end{process}

\begin{model} \hlt{Toy SOR Allocation Model}\\
For $K$ venues, let $r_k$ be the fraction routed to venue $k$, $V(\omega)$ the child size, and $D_k(\omega)$ the executable liquidity under liquidity state $\omega$. Captured liquidity is $C(r,\omega)=\sum\limits_{k=1}^{K}\min\{r_kV(\omega),D_k(\omega)\}$, $\sum\limits_{k=1}^{K}r_k=1$.\\
An SOR learns $r=(r_1,\ldots,r_K)$ by increasing allocation to venues that fill reliably and reducing allocation to venues that disappoint. The $\min$ term is the key intuition: a router cannot capture more than it sends to a venue or more than the venue actually has available.
\end{model}

\begin{mechanism} \hlt{SOR Share Feedback}\\
If $p_k(t)$ is venue $k$'s market share and $a$ is the fraction of participants using SORs, a simple share update is $p_k(t+1)=(1-a)p_k(t)+a r_k(t)$. SORs therefore turn routing weights into market share; venues with reliable liquidity gain flow, while noisy venues lose allocation slowly. When $a$ is small, historical share anchors the system; when $a$ is large, routing quality changes share quickly.
\end{mechanism}

\begin{taxonomy} \hlt{SOR Fragmentation Drivers}
\begin{enumerate}[label=\roman*.]
\setlength{\itemsep}{0pt}
\item \hlt{initial share}: incumbent venues start with the observed flow prior $p(0)$;
\item \hlt{liquidity quality}: stable depth is easier for routers to learn and trust;
\item \hlt{liquidity variance}: noisy fills push routers toward diversification;
\item \hlt{SOR adoption}: higher $a$ makes market share follow routing weights faster;
\item \hlt{reassessment speed}: faster learning accelerates convergence or instability.
\end{enumerate}
\end{taxonomy}

\begin{taxonomy} \hlt{Venue Selection Segmentation}
\begin{enumerate}[label=\roman*.]
\setlength{\itemsep}{0pt}
\item \hlt{price and depth}: displayed best quotes, hidden liquidity, and available queue size;
\item \hlt{execution probability}: fill rate, cancellation rate, queue priority, and adverse selection;
\item \hlt{economics}: commissions, access fees, maker-taker rebates, and clearing costs;
\item \hlt{access}: latency, connectivity, order types, and venue-specific restrictions;
\item \hlt{information}: pre-trade visibility, post-trade reporting, and leakage from multi-venue signalling.
\end{enumerate}
\end{taxonomy}

\begin{definition} \hlt{Market Share as Routing Prior}\\
Market share is a rough prior for where trading occurs. For stock $k$ and venue $n$, value share over $[t_1,t_2]$ is
\begin{align}
M_k(n)= \left( \sum\limits_{t_1\leq \tau_\ell^k\leq t_2}P_\ell^kV_\ell^k\mathbf{1}_n(\delta_\ell^k) \right) / \left( \sum\limits_{t_1\leq \tau_\ell^k\leq t_2}P_\ell^kV_\ell^k \right) \nonumber
\end{align}
Here $M_k(n)$ is stock $k$'s value share on venue $n$, $\ell$ indexes trades, $\tau_\ell^k$ is trade time, $P_\ell^k$ is price, $V_\ell^k$ is size, $\delta_\ell^k$ is venue, and $\mathbf{1}_n(\delta_\ell^k)$ equals $1$ if trade $\ell$ is on venue $n$ and $0$ otherwise.\\
Trade-count share replaces traded value by $\mathbf{1}_n(\delta_\ell^k)$ and divides by the number of trades. Both measures should separate continuous trading from fixing auctions.
\end{definition}

\begin{caveat} \hlt{Metric Choice}\\
The correct fragmentation metric depends on the question. Use all-in market share for venue revenue; use continuous-session trade share for the probability that an order executes on a venue; use value share when the unit of interest is currency traded.
\end{caveat}

\begin{taxonomy} \hlt{Market Share Segmentation}
\begin{enumerate}[label=\roman*.]
\setlength{\itemsep}{0pt}
\item \hlt{continuous ex-fixing}: best for venue choice and execution probability;
\item \hlt{opening, closing, and intraday calls}: best for reference-price and auction-liquidity questions;
\item \hlt{traded value}: best for economic exposure and venue revenue;
\item \hlt{trade count}: best for fill-frequency and small-order routing questions;
\item \hlt{stock or sector buckets}: needed because liquidity and capitalisation affect fragmentation.
\end{enumerate}
Combining these segments can make a primary venue look more competitive because it owns the fixing auctions.
\end{taxonomy}

\begin{empiricalfinding} \hlt{Intraday Venue-Share Shape}\\
Primary markets often regain share near the open and close because auctions, reference-price demand, and early order-book formation concentrate flow. MTF shares can be flatter, rising, or inverted depending on SOR adoption, fee agreements, and the presence of professional liquidity providers.
\end{empiricalfinding}

\begin{caveat} \hlt{Market Share Is Conditional}\\
Intraday market share is not a stable venue attribute. It changes with time of day, expiry windows, benchmark prices, news, and whether investors must trade on the venue that defines the relevant reference price.
\end{caveat}

\begin{definition} \hlt{Fragmentation Efficiency Index}\\
Let $q_n$ be the share of trading or liquidity observed in venue $n$, with $\sum_{n=1}^{N}q_n=1$.\\
The Fragmentation Efficiency Index (FEI) is normalised entropy: $\mathfrak{F}(q_1,\ldots,q_N)=\frac{-\sum_{n=1}^{N}q_n\log q_n}{\log N}$.\\
Here $N$ is the number of venues and $\mathfrak{F}$ is the dispersion score for the share vector $q$.\\
$\mathfrak{F}=0$ when all flow is in one venue and $\mathfrak{F}=1$ when flow is evenly split across venues.
\end{definition}

\begin{taxonomy} \hlt{Fragmentation Measurement Segmentation}
\begin{enumerate}[label=\roman*.]
\setlength{\itemsep}{0pt}
\item \hlt{trade-count FEI}: probability that the next print occurs away from the primary venue;
\item \hlt{value-share FEI}: economic dispersion of traded notional;
\item \hlt{displayed-liquidity FEI}: dispersion of visible book quantity, usually by price level;
\item \hlt{session FEI}: separate continuous trading from auctions and benchmark windows;
\item \hlt{universe FEI}: compare large caps, small caps, sectors, or index constituents separately.
\end{enumerate}
The same formula answers different questions only after the underlying share definition is fixed.
\end{taxonomy}

\begin{caveat} \hlt{FEI Venue Count}\\
FEI depends on the considered venue set. If $t$ is the measurement date and $N$ is the venue count, adding an empty venue mechanically lowers the index from $\mathfrak{F}(t)$ to $\frac{\log N}{\log(N+1)}\mathfrak{F}(t)$. Compare FEI through time only after controlling for venue openings, closures, and data coverage.
\end{caveat}

\begin{caveat} \hlt{Fragmentation of What}\\
Market-share and entropy metrics describe how traded value or trades are dispersed. They do not by themselves measure protocol costs, connectivity costs, hidden liquidity, or the difficulty of accessing each pool.
\end{caveat}

\begin{mechanism} \hlt{Aggregated Liquidity}\\
Consolidating visible depth across venues can reduce impact. A market order that would walk the book on one venue may execute at the same or better price when routed across several books. The benefit comes from replacing one steep book with several shallower local walks.
\end{mechanism}

\begin{caveat} \hlt{Intraday Non-Stationarity}\\
Venue share and traded volume are time-of-day conditional.\\
Opens, closes, lunch periods, linked-market openings, scheduled news, and expiry dates can create repeatable volume spikes, so a single daily average is a weak routing prior.
\end{caveat}

\begin{taxonomy} \hlt{Volume-Curve Segmentation}
\begin{enumerate}[label=\roman*.]
\setlength{\itemsep}{0pt}
\item \hlt{session}: open, continuous core, lunch break, close, and after-hours;
\item \hlt{event}: macro releases, earnings, index rebalances, and derivative expiries;
\item \hlt{reference price}: auction close, VWAP window, snapshot median, or settlement formula;
\item \hlt{calendar}: weekday, month-end, daylight-saving shifts, holidays, and local trading hours;
\item \hlt{venue}: primary auctions, lit continuous books, dark pools, and internalised flow.
\end{enumerate}
\end{taxonomy}

\begin{caveat} \hlt{Average Trade Size}\\
Falling average trade size does not mechanically mean smaller investor orders.\\
Fragmentation can split one aggressive order across several passive orders and venues, so trade prints reveal book mechanics as well as investor size.
\end{caveat}

\begin{empiricalfinding} \hlt{Liquidity and Capitalisation Segmentation}\\
Across stocks, fragmentation tends to increase with liquidity and market capitalisation.\\
For a given stock, high-volume or stressed days can reduce fragmentation as flow returns to the historical venue for depth, certainty, or familiar trading processes.
\end{empiricalfinding}

\begin{mechanism} \hlt{Why Large Caps Fragment First}\\
HFTs and liquidity arbitrageurs warehouse large-cap risk more easily: stock borrow is deeper, index or futures hedges are cleaner, inventory can be offset in correlated instruments. Small caps therefore catch up more slowly.
\end{mechanism}

\begin{marketrule} \hlt{Maker-Taker Pricing}\\
Maker-taker venues charge liquidity takers and pay rebates to liquidity makers. The displayed price can therefore differ from the economic price after fees and rebates.
\end{marketrule}

\begin{caveat} \hlt{Smart Routers vs Smart Fee Savers}\\
Routing cannot ignore fee schedules. Clearing and settlement may be netted by member inventory, while venue fees can depend on monthly turnover, cross-venue volume, or maker-taker status. The marginal cost of one child order is therefore path-dependent.
\end{caveat}

\begin{mechanism} \hlt{Rebate Effects}\\
Rebates attract passive flow and can deepen displayed queues, but they also change routing incentives. A broker or router must compare net execution quality, not only the visible quote or the nominal likelihood of a fill. The economic best venue is the one with the best expected all-in fill, not necessarily the best displayed price.
\end{mechanism}

\begin{caveat} \hlt{Subsidised Share}\\
Maker agreements and fee holidays can create market share before a venue has durable natural flow. When the subsidy ends, share can fall quickly unless the venue has become part of traders' routing logic.
\end{caveat}

\begin{caveat} \hlt{Post-Trade Barriers}\\
Fragmentation also depends on clearing, settlement, membership, and registration rules. A venue may fail to gain market share if post-trade access makes trading operationally expensive.
\end{caveat}

\begin{casestudy} \hlt{Spain as Fragmentation Control Case}\\
Spain's delayed MTF entry after MiFID makes it useful for separating fragmentation effects from broader European market stress. When alternative venues could not clear and settle efficiently, primary-market dominance persisted even while fragmentation advanced elsewhere.
\end{casestudy}

\begin{mechanism} \hlt{Event-Driven Volume Clustering}\\
Volume clusters when many traders need to act at the same reference time.\\
Linked market openings import information; scheduled macro releases create common reassessment points; derivative expiries force cash-future convergence and settlement trades.
\end{mechanism}

\begin{caveat} \hlt{Reference Rules Create Volume}\\
If a benchmark is computed from a close, auction price, or short snapshot window, traders with benchmark risk are pulled toward that window. Volume then reflects the rule as much as investor preference.
\end{caveat}

\begin{caveat} \hlt{Duplicate Liquidity}\\
The same HFT market maker may quote on several venues and cancel remaining quotes once one venue trades. Aggregated displayed depth may overstate executable liquidity unless router discounts correlated cancellations.
\end{caveat}

\begin{mechanism} \hlt{Spread-Share Feedback}\\
Best-execution and SOR logic move flow toward venues with better net spreads and reliable fills.\\
More flow attracts market makers, which can tighten spreads further. The same loop can run backward when spreads deteriorate or fills become volatile. This is why venue quality can become self-reinforcing after a venue crosses a minimum flow threshold.
\end{mechanism}

\begin{taxonomy} \hlt{Venue Quality Segmentation}
\begin{enumerate}[label=\roman*.]
\setlength{\itemsep}{0pt}
\item \hlt{tightness}: quoted and effective bid-ask spread;
\item \hlt{depth}: displayed size and replenishment near the best limits;
\item \hlt{resilience}: continuity of liquidity during volatile intervals;
\item \hlt{coverage}: how evenly liquidity is supplied across index constituents;
\item \hlt{access}: members, connectivity, latency, order types, and data quality.
\end{enumerate}
\end{taxonomy}

\begin{definition} \hlt{Coverage Liquidity Metric}\\
Coverage asks how many stocks are needed to explain a target share of index turnover, commonly 80\%.\\
Low coverage means liquidity is concentrated in a few names; high coverage means liquidity is more balanced across constituents.
\end{definition}

\begin{definition} \hlt{Gini Liquidity Imbalance}\\
A turnover Gini index summarises inequality in trading activity across stocks. A value near zero means turnover is evenly distributed; a value near one means a small subset of stocks dominates index liquidity.
\end{definition}

\begin{mechanism} \hlt{HFT Perimeter Extension}\\
Fast market makers usually enter liquid large caps first because hedging, borrowing, and inventory recycling are easier. As competition compresses returns, they extend toward less liquid names, improving coverage until hedging cost and inventory risk become uneconomic.
\end{mechanism}

\begin{caveat} \hlt{Primary Outage Dependency}\\
Many alternative venues peg orders, route decisions, or price controls to primary-market quotes. In a primary outage, those may pause alternative trading or create separate prices on venues that continue to trade.
\end{caveat}

\begin{principle} \hlt{Fragmentation Needs Shared References}\\
Fragmented markets work best when participants share timely quotes, trades, and reference prices. Competition adds liquidity only when the information layer keeps venues synchronised.
\end{principle}

\begin{caveat} \hlt{Transparency Is Not Certainty}\\
Pre-trade transparency reveals quotes, not guaranteed liquidity.\\
Post-trade transparency gives more reliable prints after the fact, so small dark pools with weak pre-trade disclosure may still be informative through execution records.
\end{caveat}

\begin{practitionernote} \hlt{Fragmented Routing Objective}\\
A practical router ranks venues by expected net cost:\\
$\text{expected cost}
= \text{spread/impact}+\text{fees}-\text{rebates}+\text{latency risk}+\text{information cost}$.\\
Patient orders can seek rebates and hidden liquidity; urgent orders pay for immediacy and certainty.
\end{practitionernote}

\begin{empiricalfinding} \hlt{Primary-Market Variance Feedback}\\
If an alternative venue delivers volatile fills, rational routers reduce future allocation to it.\\
Lower variance makes routing fractions easier to estimate; higher variance pushes flow back to the primary market and can lower the alternative venue's share.
\end{empiricalfinding}
